\chapter{Testing research code}
\label{ch:testing}

\section{Why test}\label{why-test}

A script that runs without errors can still be wrong. A test runs your code and checks the answer against something you know.

\textbf{calc.py}

\lstinputlisting{../static/examples/calc.py}

\textbf{test\_calc.py}

\lstinputlisting{../static/examples/test_calc.py}

\begin{lstlisting}[language=bash]
pip install pytest
python -m pytest
\end{lstlisting}

\textbf{What you should see}

\begin{lstlisting}[language=text]
$ python -m pytest
============================= test session starts =============================
platform win32 -- Python 3.12.12, pytest-9.2.0, pluggy-1.7.0
rootdir: ...\py
plugins: platformdirs-4.14.2
collected 3 items

test_calc.py ...                                                         [100%]

============================== 3 passed in 0.09s ==============================
\end{lstlisting}

Each dot is a passing test. To see a failure, I changed \texttt{u\ *\ area\ *\ delta\_t} to \texttt{u\ *\ area\ +\ delta\_t} in \texttt{calc.py}, as if by a slip of the finger. The code still runs without any error, but the tests catch it:

\textbf{What you should see when a test fails (first lines trimmed)}

\begin{lstlisting}[language=text]
test_calc.py FF.                                                         [100%]

================================== FAILURES ===================================
______________________________ test_known_answer ______________________________

    def test_known_answer():
>       assert heat_loss(2.1, 10, 15) == pytest.approx(315)
E       assert 36.0 == 315 ± 3.1e-04
E
E         comparison failed
E         Obtained: 36.0
E         Expected: 315 ± 3.1e-04

test_calc.py:8: AssertionError
________________ test_no_temperature_difference_means_no_loss _________________

    def test_no_temperature_difference_means_no_loss():
>       assert heat_loss(2.1, 10, 0) == 0
E       assert 21.0 == 0
E        +  where 21.0 = heat_loss(2.1, 10, 0)

test_calc.py:12: AssertionError
=========================== short test summary info ===========================
FAILED test_calc.py::test_known_answer - assert 36.0 == 315 ± 3.1e-04
FAILED test_calc.py::test_no_temperature_difference_means_no_loss - assert 21...
========================= 2 failed, 1 passed in 0.33s =========================
\end{lstlisting}

\section{Kinds of test for research}\label{kinds-of-test-for-research}

{\def\LTcaptype{none} % do not increment counter
\begin{longtable}[]{@{}
  >{\raggedright\arraybackslash}p{(\linewidth - 4\tabcolsep) * \real{0.3333}}
  >{\raggedright\arraybackslash}p{(\linewidth - 4\tabcolsep) * \real{0.3333}}
  >{\raggedright\arraybackslash}p{(\linewidth - 4\tabcolsep) * \real{0.3333}}@{}}
\toprule\noalign{}
\begin{minipage}[b]{\linewidth}\raggedright
Test
\end{minipage} & \begin{minipage}[b]{\linewidth}\raggedright
Question
\end{minipage} & \begin{minipage}[b]{\linewidth}\raggedright
Example
\end{minipage} \\
\midrule\noalign{}
\endhead
\bottomrule\noalign{}
\endlastfoot
\textbf{Known answer} & Does it match a hand calculation? & 2.1 x 10 x 15 = 315 \\
\textbf{Limit case} & Does it behave at the extremes? & No temperature gap gives no heat loss \\
\textbf{Invariant} & Does a rule always hold? & Total price equals total value \\
\textbf{Simulated data} & Can it recover a truth you planted? & A model finds a jump you added \\
\textbf{Data check} & Is the data as expected? & No duplicates, values in range \\
\textbf{Independent recalculation} & Do two separate methods agree? & The same estimate from different code \\
\textbf{Rendered output} & Does the result look right to readers? & A browser opens each page \\
\end{longtable}
}

\section{Lessons from real faults}\label{lessons-from-real-faults}

These all happened in my own projects:

\begin{itemize}
\tightlist
\item
  \textbf{An independent recalculation} confirmed an estimate exactly, and adding a check for correlated data showed its confidence interval was too narrow.
\item
  \textbf{A check of stated numbers against results} found a sample size typed wrongly into a paper.
\item
  \textbf{A test for labels} found age groups named for the wrong years, though the numbers ran fine.
\item
  \textbf{A browser test} found broken links, unreadable buttons and maths that displayed as raw code. Its first version missed some faults, because it looked for error messages and these made none.
\end{itemize}

So a passing test only shows that the things you tested held. When a fault gets through, add a test for that kind of fault.

\section{What tests cannot do}\label{what-tests-cannot-do}

They can\textquotesingle t show that a method is valid, or an argument sound. That takes reasoning and readers.

